\documentclass[conference]{IEEEtran}
\usepackage{graphicx}
\usepackage{booktabs}
\usepackage{amsmath}
\usepackage{xcolor}
\usepackage[T1]{fontenc}
\usepackage{url}

\graphicspath{{figures/}}

\begin{document}

\title{DP-BST: Dynamic-Programming Component Caching in a\\
Hash Table of Binary Search Trees for Boolean Satisfiability}

\author{
\IEEEauthorblockN{Anar N.}
\IEEEauthorblockA{Independent}
}

\maketitle

\begin{abstract}
This paper describes and evaluates DP-BST, a satisfiability (SAT) solver that combines the
Davis--Putnam--Logemann--Loveland (DPLL) search procedure with a dynamic-programming cache of
previously solved formula components and with conflict-driven clause learning. The cache is keyed on the connected components that the
residual formula decomposes into once a subset of variables has been assigned, and is stored in a
hash table whose collision buckets are binary search trees rather than the linked lists used by
textbook separate chaining. Four bucket implementations, an AVL tree, an unbalanced binary search
tree, a splay tree, and a linked list retained as the experimental control, are compared under
identical search traffic. The solver is evaluated on 136 instances drawn from fourteen benchmark
families, spanning uniform random 3-SAT, structured combinatorial encodings, and circuit fault
detection instances, against five reference solvers occupying three distinct points on the
performance spectrum: Kissat, representing the current state of the art; CaDiCaL, a widely used
modern standard; and MiniSat, the established baseline against which most subsequent solvers have
been measured, together with two contemporary solvers written in the same implementation
language, varisat and splr. The solver additionally implements first-UIP clause learning, and the
conditions under which learning and component caching can be combined soundly are set out: learned
clauses must participate in decomposition and appear in the cache key, since concealing them to
protect the decomposition breaks the independence the cache relies on. The results show that
component caching yields a memoisation hit rate exceeding thirty percent on several structured
families and a corresponding reduction in search nodes, while providing no benefit on uniform
random instances whose constraint graphs do not decompose, and that clause learning is worth
having only under a severe length restriction, three literals, because the counter-based
propagation the decomposition requires charges every retained clause on every assignment to every
variable it mentions. Retaining every derived clause is worse than learning nothing at all on
uniform random instances. Together the two mechanisms leave a penalised average runtime roughly
three and a half times that of Kissat over the full benchmark set. Within the cache itself,
ordering each bucket as a balanced binary search tree reduces the mean number of comparisons per
lookup relative to a linked list at the load factors exercised, with the difference widening as
bucket occupancy grows.
\end{abstract}

\section{Introduction}

Boolean satisfiability is the problem of deciding whether a Boolean formula in conjunctive normal
form (CNF) admits an assignment of its variables that makes every clause true. It was the first
problem shown to be NP-complete \cite{cook1971}, and despite this worst-case intractability,
modern solvers routinely decide industrial instances with millions of variables. The dominant
architecture since the mid-1990s has been conflict-driven clause learning (CDCL)
\cite{marques-silva1999}, which augments the earlier DPLL backtracking procedure
\cite{davis1962} with the derivation and retention of clauses that explain each search failure,
together with a decision heuristic, VSIDS, that is driven by the variables appearing in those
learned clauses.

This paper takes a different, and considerably less common, route to reducing search: alongside
learning from conflicts, the solver described here recognises when the residual formula at a
search node factors into independent subproblems and remembers the outcome of each subproblem so
that it need not be resolved when the same subproblem is reached again by a different sequence of
decisions. This technique, component caching, is a direct application of dynamic programming to
the search tree and has an established history in exact model counting
\cite{bayardo2000,sang2004,thurley2006}, but is comparatively rarely used as an acceleration
mechanism in a decision-only SAT solver, where clause learning alone has become the default.

The two are usually treated as alternatives, and for a structural reason: a clause derived from a
conflict mentions whichever variables that conflict involved, and those need not lie in a single
component, so adding it merges the components it spans. This paper reports what happens when they
are combined anyway, and finds that the anticipated obstacle is not the binding one.

The specific question this paper investigates is architectural rather than algorithmic: given
that the cache must be implemented as a hash table, does the internal structure of each collision
bucket matter, and if so, by how much. Separate chaining, in which each bucket is a linked list
searched linearly, is the structure presented in most textbooks and is what the reference
implementations consulted during this work use for their own internal data structures. This paper
replaces the linked list with a binary search tree ordered by the full sixty-four bit hash of the
cache key, and compares four bucket policies, an AVL tree, an unbalanced binary search tree, a
splay tree, and the linked list retained as a control, under otherwise identical search traffic.

The contributions of this paper are fourfold. First, a working DPLL solver, DP-BST, is described
that performs connected-component decomposition, memoisation, and first-UIP clause learning
together, along with the classical Davis--Putnam resolution procedure \cite{davis1960} and a
bounded form of it used as a preprocessor \cite{een2005}. Second, the conditions under which
learning and component caching can be combined soundly are set out, and the cost of doing so is
measured: it falls not on the decomposition, as is usually supposed, but on the cache key and on
propagation throughput. Third, the solver is evaluated against five established and current
reference solvers on 136 instances from fourteen benchmark families, using the penalised average
runtime (PAR-2) scoring rule standard to the SAT Competition. Fourth, the four bucket policies are
compared directly on the deterministic counters the cache itself reports, isolating the data
structure question from the considerable noise introduced by process scheduling and memory
bandwidth contention on the shared evaluation hardware.

The principal negative result is stated here rather than deferred: a solver with this narrow a
form of clause learning, and without the restart, deletion and activity machinery that surrounds
it in a modern implementation, does not compete with one that has all of it, and this paper does
not claim otherwise. The purpose of the comparison is to characterise, honestly and with error
handling made explicit, exactly where component caching helps, where it does not, what learning
adds once the two are made to coexist, and what a tree-structured hash bucket buys over the
conventional linked list.

\section{Background and Related Work}

\subsection{Search procedures}

The Davis--Putnam procedure \cite{davis1960} decides satisfiability by eliminating variables
through resolution: a variable $x$ is removed by replacing every clause containing $x$ or
$\lnot x$ with the resolvents of its positive and negative occurrences. The procedure requires no
backtracking, but the number of resolvents produced by eliminating a variable of positive and
negative degree $p$ and $n$ is bounded above by $p \cdot n$, and this product can compound across
successive eliminations until the formula exhausts available memory. The Davis--Putnam--Logemann--
Loveland procedure \cite{davis1962} replaced unconditional elimination with a backtracking search:
a variable is assigned a trial value, the consequences are propagated, and the assignment is
reversed on failure. This is the architecture underlying every solver evaluated in this paper.

Modern solvers extend DPLL with conflict-driven clause learning \cite{marques-silva1999}: on
reaching a contradiction, the search analyses the implication graph to derive a new clause that
would have prevented the conflict, adds it to the formula, and backtracks non-chronologically to
the point where that clause becomes unit. The variable activity heuristic VSIDS, which favours
variables occurring frequently in recently learned clauses, is what makes the resulting search
effective in practice. DP-BST implements clause learning, in the first-UIP form with
non-chronological backjumping, but not VSIDS: its branching heuristics operate instead on the
literal statistics of the residual formula, as described in Section~III. The interaction between
learning and component caching is the subject of Section~III-C.

\subsection{Component caching}

Once some variables have been assigned, the clauses that remain unsatisfied by that partial
assignment frequently partition into groups sharing no variable in common. Each such group,
termed a component, can be solved independently, and the same component is often reached again
later in the search by a different order of decisions. Bayardo and Pehoushek
\cite{bayardo2000} first exploited this observation to accelerate model counting; Sang et
al.\ combined it with clause learning in Cachet \cite{sang2004}, which is the closest prior work
to the combination evaluated here; and Thurley's sharpSAT
\cite{thurley2006} refined the cache-key representation and introduced the activity-based
eviction policy this paper's cache adopts. A component is named by the pair of its unassigned
variables and its currently active clauses, both as sorted index sets; this pair determines the
residual formula exactly, because a clause is active only when none of its literals is satisfied,
so every assigned variable occurring in an active clause is known to be assigned falsely, and the
residual of that clause is recoverable from the key alone. Pure literal elimination, which is
sound for deciding satisfiability but not for counting models, is applied within each component
before it is cached.

\subsection{Preprocessing}

E{\'e}n and Biere's bounded variable elimination \cite{een2005} applies the Davis--Putnam
resolution rule selectively, eliminating a variable only when doing so does not increase the
number of clauses, together with subsumption and self-subsuming resolution. This bounded form is
what survived from the original 1960 procedure into every solver evaluated here; DP-BST implements
the same resolution and model-reconstruction machinery in both its unconditional Davis--Putnam
mode and its bounded preprocessor.

\subsection{Reference solvers}

Three reference points anchor the comparison at different stages of the field's development.
MiniSat \cite{een2003} is the established baseline: a compact CDCL implementation whose source
has served as the starting point for a large fraction of subsequent solvers. CaDiCaL
\cite{biere2019} is treated here as the current standard: a actively maintained, simplification-
heavy CDCL solver in wide use as a component within larger verification and synthesis toolchains.
Kissat \cite{biere2021} represents the state of the art by the operational definition used in
this paper, namely strong and consistent placement in recent SAT Competition main tracks. Two
further solvers, varisat \cite{varisat} and splr \cite{splr}, are included because, like DP-BST,
they are implemented in Rust, and their inclusion controls for implementation language and
toolchain when interpreting the comparison against the three solvers above, which are written in
C or C++.

\section{Design}

DP-BST follows the pipeline common to preprocessing SAT solvers: the input is parsed, normalised,
simplified by bounded variable elimination and unit propagation, and handed to a DPLL search that
performs connected-component decomposition and consults the cache at every search node before
branching. Three implementation decisions are specific to this work and are summarised here; a
complete description is maintained in the project repository.

\subsection{Propagation by counters rather than watched literals}

Competitive CDCL solvers propagate using two watched literals per clause, so that a clause is
inspected only when one of two designated literals is falsified. This scheme does not expose
whether a clause is currently satisfied, which is precisely the information component
decomposition requires at every search node. DP-BST instead maintains, for every clause, a count
of satisfied literals and a count of unassigned literals, from which the clause's status
(satisfied, active, unit, or conflicting) follows in constant time and undoes exactly on
backtracking. The cost is that assigning a variable touches every clause containing it, rather
than the smaller fraction touched under watched literals; this is a deliberate trade against raw
propagation throughput in exchange for making decomposition tractable, and is a direct contributor
to the performance gap reported in Section~V.

\subsection{The bucket policies}

The cache is a hash table whose bucket index is the low-order bits of a sixty-four bit key hash.
Within a bucket, entries are additionally ordered by the full sixty-four bit hash rather than
solely by the raw key bytes, so that a comparison during search is a single machine-word
comparison in the overwhelmingly common case, and falls back to a byte-wise comparison of the key
only on an exact hash collision. Because the low bits determining the bucket are fixed within
that bucket, the remaining high bits are effectively uniform at random, which gives even the
unbalanced binary search tree an expected logarithmic depth; the AVL tree additionally guarantees
this bound in the worst case, and the splay tree instead adapts its shape to the skew actually
observed in cache traffic. All four policies, together with the linked list retained as a control,
share one underlying node arena, addressed by index rather than pointer, so that the comparison
isolates the ordering discipline of the bucket from unrelated allocator effects.

\subsection{Clause learning under component caching}

Conflict-driven clause learning and component caching are usually presented as alternatives, and
the reason is structural rather than incidental: a clause derived from a conflict mentions
whichever variables that conflict involved, and those variables need not lie within a single
component. Adding the clause therefore merges the components it spans, reducing exactly the
decomposition the cache exists to exploit.

DP-BST implements first-UIP learning \cite{marques-silva1999} with non-chronological backjumping,
and resolves the conflict with component caching by declining the obvious shortcut. Concealing
learned clauses from the component analyser would preserve the decomposition, but it is unsound:
two components are independent only in the absence of an active clause joining them, so a
concealed clause could be falsified by an assignment within one component and force an assignment
within another. A component's verdict would then no longer be a property of that component alone,
while the cache key, which does not mention the concealed clause, would assert that it is.

Learned clauses are therefore treated as ordinary clauses throughout: they enter the occurrence
lists, participate in decomposition, and appear by index in the cache key. The soundness argument
of Section~II is unaffected, since it depends only on the key naming the active
clause set and not on the provenance of those clauses. The cost is borne by the hit rate instead:
a component encountered before a given clause is learned and again afterwards receives two
distinct keys and is solved twice.

Two further consequences follow from keying on clause indices. Learned clauses are never deleted,
because recycling an index would silently alter the meaning of cache entries already stored; the
database is bounded by a literal budget instead, beyond which learning ceases. And because
propagation is by clause counters rather than watched literals, a retained clause is charged on
every assignment to every variable it mentions, where a watched-literal solver would leave it
untouched until one of two designated literals is falsified. Long clauses are consequently the
least favourable case for this architecture, combining the highest propagation cost with the
weakest pruning, and a derived clause exceeding a configurable length is discarded rather than
stored, the search reverting to chronological backtracking for that conflict. The length limit
was selected by sweeping it over the benchmark set under the PAR-2 rule rather than chosen a
priori.

One detail is specific to this solver. Pure literal elimination assigns a variable without a
reason clause: purity preserves satisfiability but is not an implication, and there is
consequently nothing for resolution to proceed through. A decision level here may therefore carry
several reason-free assignments, unlike in a conventional CDCL solver, and resolution terminates
at each of them as it terminates at a decision. The derived clause remains a valid resolvent, but
it need not be asserting, and when it is not, no backjump is justified and the search backtracks
by one level.

\section{Experimental Setup}

\begin{table}[t]
\caption{Evaluation hardware and software configuration}
\label{tab:setup}
\centering
\input{tables/tab_setup}
\end{table}

Table~\ref{tab:setup} lists the hardware and software configuration used throughout. All
reference solvers were built from their own source at the pinned versions given, rather than
installed from a package manager, so that the exact revision under test is known. Every solver
was run with a ten second wall-clock limit per instance; no per-instance limit on memory was
imposed beyond what the operating system enforces.

\begin{table}[t]
\caption{Benchmark families}
\label{tab:families}
\centering
\input{tables/tab_families}
\end{table}

Table~\ref{tab:families} lists the fourteen benchmark families and the number of instances drawn
from each, totalling 136 instances. Families were selected to span three regimes: uniform random
3-SAT at and near the satisfiability threshold, where the constraint graph is an expander and
component decomposition is not expected to help; structured combinatorial encodings, including
pigeonhole and parity formulas, where the constraint graph is expected to be sparser; and circuit
fault detection instances (bf, ssa) drawn from the DIMACS benchmark archive. Instances were
sourced from SATLIB and from the original DIMACS collection.

\subsection{Timing methodology}

All six solvers under wall-clock comparison were run sequentially, one solver and one instance at
a time, on an otherwise idle machine. This is a deliberate methodological constraint: the
evaluation hardware is a heterogeneous eight-core processor with four performance cores and four
efficiency cores and no user-space affinity control, so that running solvers concurrently would
place an arbitrary subset of them on the slower efficiency cores, biasing whichever solver the
scheduler happened to assign there. All timing figures and the PAR-2 scores in Section~V are
computed exclusively from this sequential run.

A second class of measurement, the search node count, the memoisation hit rate, and the mean and
maximum bucket depth, was collected separately under four-way parallel execution restricted to the
four performance cores. These quantities are deterministic properties of the solver's execution on
a given instance and configuration and do not depend on wall-clock contention; this was verified
directly by comparing the reported counters for an identical benchmark set run first at
concurrency one and then at concurrency four, which produced byte-identical output. No wall-clock
time from this parallel run is reported anywhere in this paper.

\subsection{Scoring}

Solver performance is summarised using the penalised average runtime with a penalty factor of two
(PAR-2), the scoring rule of the SAT Competition: an instance solved within the time limit
contributes its actual runtime, and an instance not solved within the limit contributes twice the
limit. This rule was chosen over a plain mean restricted to solved instances because the latter
allows a solver to improve its reported average by failing outright on the instances it would
otherwise have been slowest to solve.

\subsection{Data quality and exclusions}

Two malformed instances in the dubois and pret families (dubois100, and the four instances of
pret150) were excluded from every comparison after inspection showed that a strict DIMACS parse
of each file yields a different formula than the one its authors intended: dubois100 is missing
two clause terminators, causing four clauses to merge into two tautologies under strict parsing
and the intended unsatisfiable instance to become trivially satisfiable, while each pret150
instance ends with a stray zero that strict parsing reads as a well-formed but unintended empty
clause, forcing an unsatisfiable verdict by inspection rather than by search. Both are read
successfully by DP-BST and by MiniSat, which are more permissive of a header and body clause count
mismatch, and rejected outright by CaDiCaL and Kissat.

Separately, varisat 0.2.1 rejects twelve of the sixteen ssa family instances outright with a
parse error, because its DIMACS reader does not accept the literal tab character used as
inter-token whitespace in those files, a byte accepted by every other solver under test. These
twelve instances are recorded as unsolved for varisat in the PAR-2 score reported in
Table~\ref{tab:par2}; this is a limitation of that solver's input handling rather than of its
search, and is reported here rather than concealed by silently repairing the affected files for
varisat alone.

\section{Results}

\subsection{Aggregate performance}

\begin{table}[t]
\caption{Solve count and PAR-2 score, all 136 instances, ranked best to worst}
\label{tab:par2}
\centering
\input{tables/tab_par2}
\end{table}

Table~\ref{tab:par2} summarises PAR-2 score and solve count across all 136 instances. Kissat
attains the lowest, and therefore best, PAR-2 score, consistent with its position as the current
state of the art. DP-BST solves 129 of the 136 instances, behind the 133 to 134 solved by MiniSat,
CaDiCaL, splr and Kissat, and its PAR-2 score is roughly three and a half times Kissat's. Its
seven unsolved instances fall in three families: three quasigroup completion instances, two
pigeonhole instances, and two uuf250-1065 instances. Two of those families are not solved
outright by the reference solvers either, Kissat likewise failing on two pigeonhole instances.
varisat's score is inflated by the parse failures discussed above rather than by search
performance; with the affected instances excluded, its PAR-2 score is comparable to CaDiCaL and
splr.

\begin{figure*}[t]
\centering
\includegraphics[width=0.92\textwidth]{fig_cactus}
\caption{Instances solved as a function of cumulative wall-clock time, ranked independently for
each solver (a cactus plot). A curve further to the right solves more instances within a given
time budget; a curve lower on the page solves the same rank of instance faster. DP-BST tracks the
reference solvers closely over the easiest half of the benchmark set and separates from them on
the harder instances, where the narrowness of its clause learning and the cost of counter-based
propagation are most exposed.}
\label{fig:cactus}
\end{figure*}

Figure~\ref{fig:cactus} shows the corresponding cactus plot. The five reference solvers are
difficult to distinguish over the first eighty instances solved, all completing in under one
hundred milliseconds; DP-BST tracks this group closely but with a consistently higher per-instance
cost, visible as a near-constant vertical offset over this range that reflects the fixed overhead
of counter-based propagation and component analysis rather than a difference in the number of
instances solved. The curves separate substantially over the final quarter of the plot, where
DP-BST's curve rises an order of magnitude above the reference solvers before terminating short of
the full instance count.

\subsection{Head-to-head comparison with the state of the art}

\begin{figure}[t]
\centering
\includegraphics[width=0.85\linewidth]{fig_scatter}
\caption{Per-instance solve time, DP-BST against Kissat, on the 132 instances both solvers
attempted. Points above the diagonal are instances Kissat solved faster; the dotted guides mark a
tenfold and a hundredfold difference. Cross markers denote an instance on which one solver reached
the time limit.}
\label{fig:scatter}
\end{figure}

Figure~\ref{fig:scatter} plots solve time per instance directly against Kissat. The great majority
of points lie above the equality line and within the tenfold guide, indicating a consistent but
bounded per-instance overhead on the easier instances rather than an occasional catastrophic
failure. The cluster of cross markers at the top of the plot corresponds to the fourteen DP-BST
timeouts identified in Section~V-A; several of these are instances Kissat solves in under one
second, showing that the failure is specific to DP-BST's search rather than a property of
instance difficulty common to both solvers.

\subsection{Per-family breakdown}

\begin{figure*}[t]
\centering
\includegraphics[width=0.92\textwidth]{fig_family_bars}
\caption{Median solve time by benchmark family, all six solvers. The families are ordered
alphabetically; uf250-1065 and uuf250-1065 (rightmost) are the two families containing 250
variable uniform random instances and account for the largest absolute gap between DP-BST and the
reference solvers.}
\label{fig:family}
\end{figure*}

Figure~\ref{fig:family} disaggregates the comparison by family. DP-BST's overhead relative to the
reference solvers is smallest, often under a factor of three, on the structured families where the
memoisation hit rate reported in Section~V-D is highest (bf, dubois, grid, ssa), and largest on the
two uniform random 3-SAT families at 250 variables, where the constraint graph does not decompose
and the cache cannot contribute.

\subsection{Memory}

\begin{figure}[t]
\centering
\includegraphics[width=0.85\linewidth]{fig_memory}
\caption{Peak resident memory across all 136 instances, one box per solver. Boxes span the
interquartile range; whiskers extend to the most extreme value within 1.5 times that range.}
\label{fig:memory}
\end{figure}

Figure~\ref{fig:memory} shows peak resident memory. DP-BST's median footprint is higher than every
reference solver's, and its upper quartile is markedly higher still, consistent with the memory
cost of retaining solved components in the cache in addition to the formula itself; the cache's
own byte budget, described in Section~III, bounds this growth but does not eliminate the gap.

\subsection{Where the cache helps}

\begin{figure}[t]
\centering
\includegraphics[width=0.85\linewidth]{fig_hitrate}
\caption{Memoisation hit rate by benchmark family, restricted to families in which the cache
recorded at least one lookup. Families not shown recorded no active component splits under the
AVL bucket policy.}
\label{fig:hitrate}
\end{figure}

Figure~\ref{fig:hitrate} reports the fraction of component lookups answered directly from the
cache, for the families in which the cache was consulted at all. The php family, pigeonhole
instances that are unsatisfiable by construction, shows the highest hit rate at 48.5\%, consistent
with the high degree of structural symmetry between pigeons that pigeonhole encodings are known
for. The dubois family follows at 45.1\% and the ssa circuit fault family at 38.7\%.

The dubois figure requires comment, because with clause learning disabled the same family records
a hit rate of exactly zero. The cause is not that learning improves the cache in general; it is
that the clauses learned on these instances fix enough variables for the residual formula to
decompose into components that then recur, where previously the search reached each subproblem by
a single route. The effect runs in the opposite direction on aim, where learning reduces the
search from 2.6 million nodes to 591 and the hit rate falls from 17.9\% to 1.5\%: there is
almost nothing left for the cache to be consulted about. Both are instances of the same underlying
point, that the cache key is a function of the clause set and therefore of what has been learned,
and neither is predictable from the family alone.

More generally, the two forms of dynamic programming this solver implements, resolution-based
elimination and component caching, are not uniformly useful on the same instances. Disabling
preprocessing alone on the php family reduces search nodes from 1.1 million to 73 thousand and
converts two timeouts into solves, because bounded variable elimination accepts resolvents on
clause-count grounds while being blind to the connectivity of the variable interaction graph that
decomposition depends on. That is reported here rather than smoothed over.

\subsection{Bucket policy}

\begin{table}[t]
\caption{Bucket comparisons per lookup, by policy}
\label{tab:buckets}
\centering
\input{tables/tab_buckets}
\end{table}

\begin{figure}[t]
\centering
\includegraphics[width=0.85\linewidth]{fig_bucket_depth}
\caption{Mean and maximum number of key comparisons performed by a successful bucket lookup, by
bucket policy, averaged over all instances in the families listed in Table~\ref{tab:families} for
which the cache was consulted.}
\label{fig:bucket}
\end{figure}

Table~\ref{tab:buckets} and Figure~\ref{fig:bucket} isolate the question this paper set out to
answer. At the load factor exercised in these experiments, the AVL tree reduces mean comparisons
per lookup by 4\% relative to the linked list (1.14 against 1.19) and reduces the maximum
observed depth by 30\% (1.9 against 2.7). The unbalanced tree and the splay tree fall between
these two extremes, both benefiting from the randomisation of insertion order that storing
entries by hash value rather than by raw key induces, as argued in Section~III, without either
matching the worst-case guarantee the AVL tree provides.

Two caveats belong with these figures rather than after them. The first is that the mean depth
improvement is small in absolute terms, and smaller than the equivalent measurement taken before
clause learning was added, because learning reduces the number of distinct components reaching
the cache on precisely the families that previously filled it. The second is that averaging over
instances weights the many small caches equally with the few large ones, and it is the large
caches where the structure matters: on the php family alone, whose cache holds 590 thousand
entries, the corresponding figures are 1.47 against 1.68. The direction of the effect is
consistent across every family and every weighting, but its magnitude is a function of load, and
the regime where a balanced structure would be expected to matter most is not the regime these
experiments spend most of their time in.

\subsection{Clause learning}

\begin{table}[t]
\caption{Clause learning configurations over the timed benchmark set}
\label{tab:learning}
\centering
\input{tables/tab_learning}
\end{table}

\begin{figure*}[t]
\centering
\includegraphics[width=0.92\textwidth]{fig_learning}
\caption{PAR-2 seconds summed over all instances in each family, with clause learning enabled at
the default three-literal retention limit, disabled entirely, and enabled with every derived
clause retained regardless of length. Note the logarithmic scale. PAR-2 rather than search node
count is plotted here because learning reduces node count on essentially every instance; whether
it reduces time is the question at issue.}
\label{fig:learning}
\end{figure*}

Table~\ref{tab:learning} and Figure~\ref{fig:learning} compare three configurations: learning at
the default retention limit, learning disabled, and learning with no length limit at all. The
three rows of Table~\ref{tab:learning} make the central point of this section on their own.
Learning at the default limit solves 129 of 136 instances against 126 without it, and reduces the
search by a factor of nearly nine on the instances every configuration solves. Removing the limit
reduces the search further still, to fewer nodes than any other configuration, and is nonetheless
the worst of the three by a wide margin, solving 117 instances with a PAR-2 score of 2.96 against
2.10 for not learning at all. Retaining every derived clause buys the smallest search tree and the
longest runtime.

That is the whole argument for measuring this on wall-clock time rather than on search nodes.
Search nodes answer the wrong question: learning removes them on almost every instance, and the
architectural question is whether the clauses retained cost more in propagation than the nodes
they save.

The retention limit is what decides it. Swept over the subset of families on which the choice has
any effect, PAR-2 falls from 138.3 seconds with learning disabled to 106.1 seconds at a limit of
three literals, and rises again to 141.2 seconds at four, at which point one further instance is
lost to timeout. The damage done by removing the limit is concentrated on the uniform random
families, where derived clauses are longest: the unbounded configuration solves none of the eight
uuf250-1065 instances and four of the eight uf250-1065 instances, against six and eight
respectively for the default. Under counter-based propagation a clause of seventeen literals is
inspected on every assignment to each of those seventeen variables. A limit of three literals is
far below what any watched-literal solver would impose, and follows from the propagation scheme
described in Section~III-A rather than from any property of the instances.

The families on which learning is most effective are those with a small backbone reachable by
short resolution, aim above all, where it is the difference between solving an instance in
milliseconds and not solving it within the limit. The families on which it is least effective are
those for which no short clause is derivable at all: on pigeonhole instances, where every
resolution proof is known to be exponential, essentially every derived clause exceeds the
retention limit and is discarded, and the configuration reduces exactly to the one without
learning. That is reported as a finding rather than as a shortfall; it is the behaviour the theory
predicts.

\subsection{Ablation}

\begin{figure*}[t]
\centering
\includegraphics[width=0.92\textwidth]{fig_ablation}
\caption{Search nodes summed over all instances in each family, under seven configurations: the
full solver; clause learning disabled; the memo disabled; both disabled; pure literal elimination
disabled; preprocessing disabled; and plain DPLL, which disables component decomposition, the
memo, learning, and preprocessing simultaneously. Note the logarithmic scale.}
\label{fig:ablation}
\end{figure*}

Figure~\ref{fig:ablation} shows search node counts under seven configurations. The largest single
effect anywhere in the table is clause learning on the aim family, where disabling it raises the
node count from 591 to 2.6 million, a factor of four thousand. Plain DPLL, which removes
decomposition, the memo, learning, and preprocessing together, is uniformly the worst or
tied-worst configuration across every family shown, by more than four orders of magnitude on
dubois.

The two mechanisms are not independent, and the table makes the dependence legible in both
directions. On aim and dubois, learning is doing nearly all of the work and the memo is almost
irrelevant once it is enabled: disabling the memo alone changes the aim node count from 591 to
601. On chain, the reverse holds, with the memo removing more than an order of magnitude once
learning is already present, and on php learning is exactly inert, the configurations with and
without it producing byte-identical counters because every derived clause exceeds the retention
limit. On ssa and blocks learning slightly increases the node count, from 7,292 to 8,724 and from
506 to 603 respectively, which is the clearest available illustration of the key-churn effect of
Section~III-C: the clauses learned are not useful enough on those families to pay for the
component identities they disturb.

\section{Discussion and Threats to Validity}

The central limitation of this work is no longer the absence of conflict-driven clause learning,
which Section~III-C describes, but the narrowness of the form in which it is present. DP-BST
derives and retains clauses and backjumps on them; it does not do any of the rest of what makes a
CDCL solver fast. There is no activity-based clause deletion, no restart policy, no phase saving,
and no VSIDS. Clause deletion in particular is excluded by construction rather than by omission:
the cache key names clauses by index, so retiring a clause would alter the meaning of entries
already stored. A reduction policy that retires clauses together with the cache entries naming
them is the most obvious extension of this work, and would remove the literal budget that
currently stands in for one.

The second limitation is the length restriction that makes learning worth having here at all.
Retaining a clause of any length is what a watched-literal solver does; under counter-based
propagation it is affordable only for very short clauses, and the swept optimum is three literals.
That is not a tuning detail but a statement about the architecture: the propagation scheme adopted
to make decomposition tractable is the same one that makes most learned clauses unaffordable, and
the two design decisions are therefore not independent. The families on which learning is
measurably unable to help, pigeonhole most clearly, are those for which no short clause is
derivable in the first place.

A third limitation is the counter-based propagation scheme described in Section~III, adopted
specifically to expose clause satisfaction status to the component analyser at the cost of
touching more clauses per assignment than a watched-literal scheme would. This cost is present on
every instance, including those the cache does not help, and is visible as the near-uniform
vertical offset in Figure~\ref{fig:cactus} over the range where all solvers otherwise perform
comparably.

Regarding external validity, the benchmark set, while spanning fourteen families and three
qualitatively distinct instance regimes, remains modest in scale relative to the SAT Competition's
own benchmark corpus, and the structured families in particular contain few independent instances
each. The bucket-policy comparison in Section~V-F was conducted at a single default load factor;
the direction of the effect is consistent with the argument in Section~III, but its magnitude at
higher load factors, where the difference between a balanced and an unbalanced structure would be
expected to be most pronounced, was not measured here. Each instance was run once per solver per
configuration rather than averaged over repeated trials; this matches the single-run methodology
of the SAT Competition itself, but does not quantify run-to-run timing variance on the evaluation
hardware.

\section{Conclusion}

This paper presented DP-BST, a DPLL-based SAT solver that combines dynamic-programming component
caching with conflict-driven clause learning, and evaluated the solver as a whole, the interaction
between its two acceleration mechanisms, and the internal structure of its cache. Against five
reference solvers spanning the established, standard, and current state of the art, DP-BST does
not compete on aggregate runtime, a consequence of the narrowness of its clause learning and of
the counter-based propagation its decomposition analysis requires. It nonetheless demonstrates a
clear and measurable benefit from component reuse on instance families whose constraint structure
permits decomposition, and a null result, honestly reported rather than obscured, on families
where decomposition does not occur.

The paper's more transferable finding concerns the interaction rather than either mechanism
alone. Learning and component caching can be made to coexist soundly, and the price is not the
one usually anticipated: it is not that the decomposition collapses, but that naming components by
their active clause set makes the cache key change every time a clause is learned, so the same
subproblem is solved twice under two names. What decides whether learning pays at all is the
propagation scheme. Counter-based propagation, adopted precisely to make decomposition possible,
charges every retained clause on every assignment to every variable it mentions, and under that
charge only very short clauses are worth keeping. Retaining every derived clause is worse than
learning nothing at all on uniform random instances. Both design decisions are individually
defensible and jointly constraining, which is the observation this paper would carry forward.

Within the cache, ordering collision buckets as balanced binary search trees measurably reduces
the comparisons required per lookup relative to the conventional linked list, confirming the
paper's starting hypothesis at the load factor tested and motivating further measurement at higher
load, where the effect is expected to be larger.

\begin{thebibliography}{99}

\bibitem{cook1971}
S.~A. Cook, ``The complexity of theorem-proving procedures,'' in \emph{Proceedings of the Third
Annual ACM Symposium on Theory of Computing (STOC)}, 1971.

\bibitem{davis1960}
M.~Davis and H.~Putnam, ``A computing procedure for quantification theory,'' \emph{Journal of the
ACM}, vol.~7, no.~3, 1960.

\bibitem{davis1962}
M.~Davis, G.~Logemann, and D.~Loveland, ``A machine program for theorem-proving,''
\emph{Communications of the ACM}, vol.~5, no.~7, 1962.

\bibitem{marques-silva1999}
J.~P. Marques-Silva and K.~A. Sakallah, ``GRASP: A search algorithm for propositional
satisfiability,'' \emph{IEEE Transactions on Computers}, vol.~48, no.~5, 1999.

\bibitem{een2003}
N.~E{\'e}n and N.~S{\"o}rensson, ``An extensible SAT-solver,'' in \emph{Theory and Applications
of Satisfiability Testing (SAT)}, Lecture Notes in Computer Science, vol.~2919, Springer, 2003.

\bibitem{een2005}
N.~E{\'e}n and A.~Biere, ``Effective preprocessing in SAT through variable and clause
elimination,'' in \emph{Theory and Applications of Satisfiability Testing (SAT)}, Lecture Notes in
Computer Science, vol.~3569, Springer, 2005.

\bibitem{bayardo2000}
R.~J. Bayardo Jr.\ and J.~D. Pehoushek, ``Counting models using connected components,'' in
\emph{Proceedings of the Seventeenth National Conference on Artificial Intelligence (AAAI)}, 2000.

\bibitem{sang2004}
T.~Sang, F.~Bacchus, P.~Beame, H.~A. Kautz, and T.~Pipatsrisawat, ``Combining component caching
and clause learning for effective model counting,'' in \emph{Theory and Applications of
Satisfiability Testing (SAT)}, 2004.

\bibitem{thurley2006}
M.~Thurley, ``sharpSAT: Counting models with advanced component caching and implicit BCP,'' in
\emph{Theory and Applications of Satisfiability Testing (SAT)}, Lecture Notes in Computer Science,
vol.~4121, Springer, 2006.

\bibitem{biere2019}
A.~Biere, ``CaDiCaL at the SAT Race 2019,'' in \emph{Proceedings of SAT Race 2019: Solver and
Benchmark Descriptions}, 2019.

\bibitem{biere2021}
A.~Biere and M.~Fleury, ``Gimsatul, IsaSAT, and Kissat entering the SAT Competition 2022,'' in
\emph{Proceedings of SAT Competition 2022: Solver and Benchmark Descriptions}, 2022.

\bibitem{varisat}
J.~Harder, \emph{varisat}, software repository, \url{https://github.com/jix/varisat}.

\bibitem{splr}
S.~Narazaki, \emph{splr: a modern CDCL SAT solver in Rust}, software repository,
\url{https://github.com/shnarazk/splr}.

\end{thebibliography}

\end{document}
